\subsection{Predeclared public-repository comparison}
The following values are generated from the fixed 26 chronological targets in \texttt{cpmpy/cpmpy}. MAE evaluates the predictive median. Width is the difference between the inverse-CDF 90th and 10th percentiles, and coverage is the proportion of observed targets within those inclusive endpoints. The requested percentiles give no guarantee of future coverage.
\begin{center}
\small
\begin{tabular}{@{}lrrrrr@{}}
\toprule
Method & Targets & Mean RPS & MAE & Mean width & Coverage \\
\midrule
Recent empirical & 26 & 2.526 & 3.038 & 13.12 & 80.8\% \\
Expanding empirical & 26 & 2.275 & 3.538 & 10.15 & 92.3\% \\
Zero & 26 & 4.615 & 4.615 & 0.00 & 15.4\% \\
Previous period & 26 & 4.500 & 4.500 & 0.00 & 15.4\% \\
Poisson, prior-four mean & 26 & 2.766 & 3.731 & 5.69 & 42.3\% \\
\bottomrule
\end{tabular}
\end{center}
The recent empirical model's mean RPS is 11.03\% higher than the expanding empirical comparator's on this fixed public-repository window. The comparison reports the observed difference without a confidence interval adjusted for dependence. It leaves general superiority and performance in the user's repositories unestablished.
The recent-window model fails the predeclared requirement of at least 5\% lower mean RPS. Its point-error result leaves that failure unchanged, so the default expanding-history rule is retained.
The complete report retains target dates, outcomes and paired losses for every method. All results are retrospective reconstructions from the retained visible inventory.
